\section{Uniform subset graphs}
Let $\Omega$ have $n$ elements and let $\Gamma(n,k,r)$ join two
$k$-subsets when their intersection has size $r$, where $0\leq r<k$.
Unless an endpoint is explicitly discussed, take $n\geq2k+1$ and
$k\geq2$. Write $N=\binom nk$, $\nu=\binom kr\binom{n-k}{k-r}$ and
$\bar\nu=N-1-\nu$. A binomial coefficient outside its ordinary range
is zero. For $|x\cap y|=j$, put
\[
 i_j=\sum_a\binom ja\binom{k-j}{r-a}^{\!2}
              \binom{n-2k+j}{k-2r+a},\qquad
 P_j=i_j+[j=r],\qquad \ell_j=\nu-P_j.
\]
The natural $S_n$ action is transitive on vertices and on each ordered
pair class of fixed intersection size. Thus it is transitive on the
nonloop edges of $\Gamma(n,k,r)$. The complement generally contains
several such classes, and loops form a separate class in the reflexive
graph; the source's phrase extending edge transitivity to both associates
cannot be taken literally for this $S_n$ action. No such additional
transitivity is used in the count argument.
These are exactly the common-neighbour and scalar products in source
Proposition~3. The four-region count is proved in
\Uref{ULCD.SubsetCounts}; \Uref{ULCD.Scalar} gives
\begin{equation}\label{eq:subset-tests}
\begin{array}{rcl}
 C\text{ is RLCD}&\Longleftrightarrow&P_j=0\text{ for all }j,\\
 C\perp\bar C&\Longleftrightarrow&P_j=\nu\text{ for all }j,\\
 C\text{ is CLCD}&\Longleftrightarrow&P_j=\nu\text{ for all }j
                      \text{ and }(\nu,\bar\nu)\ne(0,0),\\
 \bar C\text{ is RLCD}&\Longleftrightarrow&\bar\nu-\nu+P_j=0
                      \text{ for all }j.
\end{array}
\end{equation}
All equalities in these tests are in $\mathbb F_p$. In the last line the
diagonal test forces $\bar\nu=0$, so $C\perp\bar C$. If also $\nu=0$,
then $C$ is RLCD; otherwise $C=\bar C^\perp=\im(I+B)$.
This supplies all three clauses of the source proposition, with its
possible complement exception replaced by the exact constant-line defect.

For the Johnson graph $r=k-1$, the only possibly nonzero $i_j$ are
$i_{k-2}=4$, $i_{k-1}+1=n-1$ and $i_k=k(n-k)$.
For the Kneser graph $r=0$, $i_j=\binom{n-2k+j}{k}$.
For $r=1$ they are
\[
 i_j=j\binom{n-2k+j}{k-1}+(k-j)^2\binom{n-2k+j}{k-2}.
\]
Together with $\ell_j=\nu-P_j$, these give every formula in source
Example~2. Only realized intersection sizes should be tested outside
the domain $n\geq2k+1$.

\begin{corollary}\label{cor:subset-classification}
In this domain, the complete conditions for the five families listed
on source page~11 are the following. A blank condition means that there
are no further primes or residues.
\begin{center}\small
\begin{tabular}{lll}\toprule
$(k,r)$&RLCD&CLCD\\\midrule
$(2,0)$&$p=2,\ n\equiv3\pmod4$&$p=2,\ n\equiv1\pmod4$\\
$(3,0)$&$p=2,\ n\equiv1\pmod4$&$p=2,\ n\equiv1\pmod4$\\
&$p=3,\ n\equiv5\pmod9$&$p=3,\ n\equiv2,8\pmod9$\\
&$p\geq5,\ n\equiv5\pmod p$&$p\geq5,\ n\equiv5\pmod p$\\
$(4,0)$&$p=2,\ n\equiv7\pmod8$&$p=2,\ n\equiv3\pmod8$\\
$(3,1)$&$p=2,\ n\equiv0\pmod4$&$p=2,\ n\equiv0,1\pmod4$\\
$(4,1)$&$p=2,\ n\equiv1\pmod4$&$p=2,\ n\equiv1\pmod8$\\
&$p=3,\ n\equiv6\pmod9$&$p=3,\ n\equiv6\pmod9$\\\bottomrule
\end{tabular}
\end{center}
\end{corollary}
\begin{proof}
These are the five cases of \Uref{ULCD.SubsetFamilies}, with precisely
the same $(n,k,r,p)$ parameters. That result substitutes the displayed
integer counts, excludes other primes by polynomial identities and
checks the remaining complete residue periods. Thus the ``only if''
parts do not rest on a finite sample of $n$.
\end{proof}

The Johnson case of source Corollary~1 is RLCD precisely for $p=2$ and
odd $n$. It is CLCD precisely when, in addition, $\binom nk$ is even:
the degree vanishes modulo two, so the constant-line test is decisive.
This makes its printed exception exact. Domain matters: for instance
$J(n,n-1)=K_n$ is RLCD when $p\mid n-1$ and CLCD when $p\nmid n-1$;
the unrestricted wording of the source corollary cannot cover these
endpoints by the same parity criterion.
For the odd graph $KG(2k+1,k)$, the count at $j=0$ is $P_0=1$ for
$k\geq2$, so it is never RLCD. For $k\geq3$, $P_1=0$, which rules out
CLCD by comparison with $P_0$. At $k=2$ the counts are
$i_0=0$, $i_1=1$, $i_2=3$, hence $P=(1,1,3)$.
Thus all $P_j=\nu$ precisely in characteristic two, and $\nu=3\ne0$.
The binary Petersen graph is CLCD, correcting the source's universal
exclusion.

For every regular graph, RLCD implies $\nu=0$ by the diagonal of
$A(A+I)=0$, and consequently implies $C\perp\bar C$. Conversely,
$C\perp\bar C$ implies RLCD exactly when $\nu=0$. The source's example
$\Gamma(13,3,1)$ over $\mathbb F_2$ has all $P_j=\nu=1$, so is CLCD
and not RLCD. This recovers the final clause and note of Corollary~1.

\subsection{Every computational finding on source page 8}
The following table evaluates \eqref{eq:subset-tests} for all sixteen
reported cases. The column $P$ means the common residue of every
$P_j$, not merely one entry; both degrees are also reduced modulo $p$.
It is a complete small integer certificate, obtained by the displayed
binomial formula, rather than an unexplained graph-code search.
\begin{center}\small
\begin{tabular}{rrrrrrrcc}\toprule
$n$&$k$&$r$&$p$&$P$&$\nu$&$\bar\nu$&RLCD&CLCD\\\midrule
9&4&1&2&0&0&1&yes&yes\\
9&4&3&2&0&0&1&yes&yes\\
11&4&0&2&1&1&0&no&yes\\
11&4&3&2&0&0&1&yes&yes\\
13&4&1&2&0&0&0&yes&no\\
13&4&3&2&0&0&0&yes&no\\
15&4&0&2&0&0&0&yes&no\\
15&4&3&2&0&0&0&yes&no\\
17&4&1&2&0&0&1&yes&yes\\
17&4&3&2&0&0&1&yes&yes\\
15&4&1&3&0&0&2&yes&yes\\
11&5&4&2&0&0&1&yes&yes\\
12&5&3&2&0&0&1&yes&yes\\
13&5&2&2&0&0&0&yes&no\\
13&5&4&2&0&0&0&yes&no\\
15&5&4&2&0&0&0&yes&no\\\bottomrule
\end{tabular}
\end{center}
When both answers are yes, the reflexive and complementary codes are
the same orthogonal dual, as asserted in the source. All its positive
findings follow, and the negative entries are now determined as well.

\subsection{Spectra and the three dimension examples}
We retain the classical integral spectrum of the Johnson association
scheme, attributed by the source to Krebs--Shaheen, Delsarte and
Brouwer--Haemers. Its multiplicities are
$m_j=\binom nj-\binom n{j-1}$ for $0\leq j\leq k$.
The Kneser and Johnson eigenvalues are
$(-1)^j\binom{n-k-j}{k-j}$ and $k(n-k)-j(n+1-j)$, respectively.
For the general relation the retained Eberlein formula is
\[
 \epsilon_j=\sum_i(-1)^i\binom ji\binom{k-j}{k-r-i}
                                 \binom{n-k-j}{k-r-i}.
\]
The reflexive and complement spectra follow by adding one and by
subtracting from $J-I$; no further classification input is used.

For $(k,r)=(3,1)$ the eigenvalues are
\[
 3\binom{n-3}{2},\quad \binom{n-4}{2}-2(n-4),\quad-2n+11,\quad3.
\]
When $n\equiv0\pmod4$, only the last two are odd. Therefore
\Uref{ULCD.Rank} gives $m_2+m_3=\binom n3-n$, proving Example~3(1).
For $(k,r)=(4,1)$ they are
\[
 4\binom{n-4}{3},\quad \binom{n-5}{3}-3\binom{n-5}{2},\quad
 -(n-6)(n-9),\quad3n-22,\quad-4.
\]
When $p=2$, $n\equiv1\pmod4$, only the fourth is nonzero, so the
dimension is $m_3=\binom n3-\binom n2$.
When $p=3$, $n\equiv6\pmod9$, exactly the fourth and fifth survive,
giving $m_3+m_4=\binom n4-\binom n2$.
The degree-$4\binom{n-4}{3}$ row is a nonzero word and supplies the two
stated distance upper bounds. These are source Example~3(2)--(3), with
no claim that either upper bound is always attained.
